% ============================================================
%  H. Høffding, Religionsfilosofi
%  Kjøbenhavn og Kristiania: Gyldendalske Boghandel
%  Nordisk Forlag, 1901. Andet Oplag 1906.
%  TRANSCRIPTION (Danish)
%  Source: KB scan (pdf) and modern epub,
%  cross-checked against each other.
% ============================================================
\documentclass[12pt,a4paper]{book}
\usepackage[utf8]{inputenc}
\usepackage[T1]{fontenc}
\usepackage[danish]{babel}
\usepackage{microtype}
\usepackage{setspace}
\usepackage[margin=1.2in]{geometry}
\usepackage{fancyhdr}
\setlength{\headheight}{14.5pt}
\addtolength{\topmargin}{-2.5pt}
\usepackage{hyperref}

\hypersetup{
  pdftitle={Religionsfilosofi},
  pdfauthor={Harald Høffding},
  hidelinks
}

\pagestyle{fancy}
\fancyhf{}
\fancyhead[LE,RO]{\thepage}
\fancyhead[RE]{\textit{Religionsfilosofi}}
\fancyhead[LO]{\textit{\leftmark}}

\setstretch{1.3}

\title{\textbf{Religionsfilosofi}}
\author{Harald Høffding}
\date{Kjøbenhavn og Kristiania:\\
  Gyldendalske Boghandel, Nordisk Forlag\\
  1901. Andet Oplag 1906.}

\begin{document}
\maketitle
\thispagestyle{empty}

\vspace*{\fill}
\begin{center}
\textit{Alte dubitat, qui altius credit.}\\[6pt]
(Dyb Tvivl, paa Grundlag af endnu dybere Tro.)
\end{center}
\vspace*{\fill}
\newpage

\tableofcontents
\newpage

%% ============================================================
%%  I. OPGAVE OG FREMGANGSMAADE
%% ============================================================
\chapter{Opgave og Fremgangsmaade}
\label{ch:opgave}

\begin{flushright}
\textit{ou lyseis agnoōn ton desmon.}\\
\textsc{Aristoteles}.\\[4pt]
\textit{(Man kan ikke løse uden at vide, hvor Baandet er).}
\end{flushright}

\bigskip

\noindent\textbf{1.} Ved Religionsfilosofi kunde man enten forstaa en
Tænkning, der udspringer af Religion og bygger paa Religion, eller en
Tænkning, der har Religion til sin Genstand. Det er i denne sidste
Betydning, Ordet her skal tages. Den Tænken, der udspringer af og
bygger paa Religion, hører selv med til Religionsfilosofiens Genstand
og bør ikke kaldes Religionsfilosofi.

Ved Religion tænke vi nærmest paa en sjælelig Tilstand, i hvilken
Følelse og Trang, Frygt og Haab, Begejstring og Hengivelse spille en
større Rolle end Eftertanke og Gransken, og i hvilken Anskuelse og
Billede herske mere end Undersøgelse og Reflexion. Vel vil der atter
og atter i det religiøse Liv opstaa en uvilkaarlig Trang til at gøre
sig Rede for sin Tilstand og for sit Indhold, --- for sine Motivers
Værd og for sine Tankers Gyldighed. Men den religiøse Tænkning, der
saaledes opstaar, stiller ikke det religiøse Problem paa afgørende
Maade. Dens Problemer opstaa indenfor Religionen, men selve
Religionen bliver ikke Problem. Der kendes endnu ikke andet
Udgangspunkt end Religionen selv; i hvert Tilfælde ere andre
Udgangspunkter saa underordnede i Forhold til selve Religionen, at de
ikke faa afgørende Betydning. Af denne Art er den religiøse Tænken i
de Tider, i hvilke Religionen er Et og Alt paa det aandelige
Omraade. Det er Religionens klassiske Tider, dels de store
Begyndelsestider, i hvilke den med Oprindelighedens hele Magt drager
al Kraft og Interesse til sig, dels de store Organisationstider, i
hvilke al forhaandenværende Kultur er formet eller bøjet til Lydighed
under de højeste religiøse Ideer. I saadanne klassiske Tider er der
en stor Enhed, eller dog en stor Harmoni paa det aandelige Livs
Omraade. En saadan klassisk Tid havde Kristendommen i den
urkristelige Menigheds Tid og senere i den store Organisationsperiode
i Middelalderens Blomstringstid. Et religiøst Problem danner der sig
først, naar andre Sider af Aandslivet --- Videnskab og Kunst, Moral
og Samfundsliv --- begynde at frigøre sig og at optræde med Fordring
paa at have selvstændig Værdi. De bedømme da Religionen ud fra deres
Synspunkter og efter deres Maalestokke, medens Religionen i sine
klassiske Tider enten slet ikke ænsede dem --- som i
Begyndelsestiderne, --- eller --- som i Organisationstiderne ---
værdsatte dem ud fra sit Synspunkt og efter sin Maalestok, saa at
Spørgsmaalet blev, hvorvidt disse to forskellige Vurderinger kunde
bringes i indbyrdes Harmoni. --- Der maa være bestemte historiske
Betingelser til Stede, forat det religiøse Problem kan blive stillet
paa afgørende Maade. Det er, om man vil, kun i ulykkelige Tider, at
der existerer et religiøst Problem. Dette Problem er jo Udtryk for en
aandelig Disharmoni: Aandslivets forskellige Elementer arbejde nu
ikke mere sammen paa saa sluttet en Maade som før; de føre i
forskellige Retninger, mellem hvilke der maaske tilsidst maa vælges.
Men da er ogsaa baade Nødvendigheden og Muligheden af en Undersøgelse
til Stede.

At Religionen gøres til Problem, vil være stødende for Mange. Men
Tænkningen maa have Ret til at undersøge Alt, naar den engang er
vaagnet, og kun den selv kan afstikke sine Grænser. Hvem ellers
skulde vel afstikke dens Grænser? Den, som ikke mærker noget
Problem, har naturligvis ingen Grund til at tænke; men han har heller
ingen Grund til at hindre Andre i at tænke. Den, der frygter for at
miste sit aandelige Tilhold, maa holde sig borte. Ingen vil berøve
den Fattige hans eneste Lam --- men saa maa den Fattige heller ikke
uden Nødvendighed komme trækkende med det paa Alfarvej og fordre, at
Færdselen skal stanses for hans Skyld. Iøvrigt viser Erfaringen, at
det snarere er Oxerne end Lammene, der højrøstet melde sig med deres
Forargelse i Tide og i Utide (især i Utide). Det er ikke saa meget
de virkelige Fattige i Aanden, som de stive og buldrende
Dogmatikere, der tage paa Vej, naar den frie Forsken kræver Ret til
at røre sig paa det religiøse Omraade som paa ethvert andet.

Den Undersøgelse, jeg her vil anstille, henvender sig hverken til de
Færdige eller til de Ængstelige. De Færdige findes i alle Lejre, ikke
mindst blandt de saakaldte ``Fritænkere'' --- en Klasse af Mennesker,
der ligesom ``Ormene'' i Linnés System kun kan karakteriseres ved
negative Prædikater, fordi den skal indbefatte saa mange forskellige
Former. De Færdige sidde inde med en bestemt negativ eller positiv
Løsning af det religiøse Problem og have derved tabt Sansen for
videre Tænkning over det. De Ængstelige tør ikke tænke over det. ---
Derimod henvender Undersøgelsen sig til de Vordende. ``En Vordende
vil altid være taknemmelig'' --- ligemeget i hvilken Retning hans
Vorden gaar.

Filosofien er jo, som bekendt, rigere paa Ideer, Synspunkter og
Drøftelser end paa definitive Resultater, og Filosofen vilde ofte
være ilde stedt, dersom han ikke troede, at den stadige
Sandhedsstræben hørte til de højeste aandelige Værdier, som det er al
sand Religions Opgave at hævde. Selvom vi ikke kunne lære Andet af en
religionsfilosofisk Undersøgelse, kunde vi dog maaske lære Noget om
Aarsagerne til den Strid, der staar om religiøse Spørgsmaal, og om
denne Strids Betydning for det aandelige Livs Udvikling. Og skulde
det religiøse Problem være uløseligt, kunde vi maaske faa at vide,
hvorfor det er uløseligt.

\medskip

\noindent\textbf{2.} Religionsfilosofien bør ikke gaa ud fra et
færdigt filosofisk System. I en vis Forstand maa al Tænken være
systematisk, ti ``System'' betyder ordret ``det, som staar sammen'',
og det er jo en ubønhørlig Fordring, der stilles til vore Tanker, at
de maa kunne ``staa sammen'' med hverandre. Der er en indre
Sammenhæng og Konsekvens, som paa alle Omraader er det sidste
Kendemærke paa Sandheden, og som derfor ogsaa maa kræves paa det
religionsfilosofiske Omraade. Men det religionsfilosofiske Arbejde
vil være frugtbarest, hvor Religionen og dens Fænomener ikke stilles
i Forhold til en allerede afsluttet filosofisk Opfattelse, men
belyses af en Filosoferen, der netop er ved at undersøge, om en
Afslutning idethele er mulig. Opgaven er at drøfte Religionens
Forhold til Aandslivet. Religion er selv en Art af eller Form for
aandeligt Liv, og den vurderes kun ret ved at ses i sit Forhold til
andre Arter eller Former. Det er ingen fremmed Maalestok, der herved
anlægges paa Religionen. Ti der er en Tjeneste, som ingen aandelig
Magt --- om den end optræder med de største Navne --- kan unddrage
sig, nemlig at tjene det aandelige Livs dybe og rige Udvikling. Den
religionsfilosofiske Opgave vil da først og fremmest gaa ud paa at
finde Midler til at undersøge, hvorvidt Religionen under vor
aandelige Kulturs Vilkaar vedblivende kan yde saadan Tjeneste. Til
Løsningen, eller bedre: Behandlingen af denne Opgave kan
Religionsfilosofien ikke møde med andre Forudsætninger end dem,
enhver anden Videnskab medbringer til Behandlingen af sit Emne. Meget
ofte vil den sikkert blive nødt til at erklære, at der er
Spørgsmaal, som ikke kunne besvares, fordi Forholdene paa det
aandelige Omraade ere mere indviklede og mere forskelligartede, end
positive og negative Dogmatikere mene. Derfor er det dog ikke uden
Betydning, at Problemerne drøftes. Den bevidste Indsigt i Løsningens
Umulighed (\textit{docta ignorantia}) opfordrer jo naturligt til at
undersøge, om det da er saa aldeles nødvendigt, at der findes en
Løsning. Desuden vil Blikket blive skærpet for de særlige Forhold,
under hvilke vor aandelige Kultur lever, og som sandsynligvis ikke
vilde vedblive at bestaa, ligesaa lidt som de allerede nu øve lige
stor Indflydelse paa alle Personligheder.

En religionsfilosofisk Undersøgelse vil hensigtsmæssigt kunne
anlægges saaledes, at den bestaar af en erkendelsesteoretisk, en
psykologisk og en etisk Del.

I sine klassiske Tider tilfredsstiller Religionen alle Menneskets
aandelige Fornødenheder, ogsaa dets Erkendelsestrang. Religiøse Ideer
indeholde da for Mennesket Forklaringen af Tilværelsen i dens Helhed
og i dens enkelte Dele. Men naar der har udviklet sig en selvstændig
Videnskab, fremtræder der en anden Art af Forstaaelse eller
Forklaring end den religiøse, og Spørgsmaalet bliver, om disse to
Arter ere forenelige. Muligvis kunde Forholdet være det, at den nye,
videnskabelige Forklaringsmaade mere og mere afløser den religiøse
Forklaring, hvad Enkelthederne angaar, men at den ikke kan anvendes
overfor Tilværelsen som Helhed betragtet, --- overfor hvad man plejer
at kalde de ``første'' eller de ``sidste'' Spørgsmaal. Da bliver det
at drøfte, hvorvidt Religionen formaar at løse de Gaader, som
Videnskaben ikke kan løse. Dersom det saa skulde vise sig, at
religiøse Forestillinger, hverken hvad det Hele eller det Enkelte
angaar, kunde hjælpe til Noget, der med Rette skulde kunne kaldes
Forstaaelse eller Forklaring, saa bliver det Spørgsmaal at opkaste,
hvilken Betydning disse Forestillinger da kunne have.

Dette sidste Spørgsmaal fører over fra den erkendelsesteoretiske
Religionsfilosofi til den psykologiske. Ti dersom de religiøse
Forestillinger ikke mere have Erkendelsesværdi, maa den Værdi, de
kunne besidde, bestaa i at være Udtryk for andre Sider af Sjælelivet
end den intellektuelle. Der maa da forsøges en Beskrivelse af det
religiøse Sjæleliv, særligt af Forestillingernes Forhold til religiøs
Erfaring og religiøs Følelse. Derigennem vil det vise sig, at
Religionen i sit inderste Væsen ikke har med Forstaaelsen af
Tilværelsen, men med Vurderingen af den at gøre, og at de religiøse
Forestillinger udtrykke det Forhold, i hvilket den faktiske
Tilværelse, saaledes som vi kende den, staar til det, som for os
giver Livet dets højeste Værdi. Religionens Kærne er --- efter den
Hypotese, vi ledes til --- en Tro paa Værdiens Bestaaen i
Tilværelsen. Denne Tro fremtræder i store, anskuelige Billeder i
Folkereligionerne, især i de højeste af disse. Ogsaa hos dem, der
staa udenfor enhver Folkereligion, vil der kunne røre sig en saadan
Tro, selvom den ikke her udformes i bestemt Skikkelse.

Overgangen fra psykologisk til etisk Religionsfilosofi motiveres
naturligt ved, at der spørges, hvilken etisk Værdi det har, at der
tros paa Værdiens Bestaaen. Etiken angaar Frembringen af Værdier, og
den Mulighed kan ikke afvises, at Arbejdet paa at frembringe Værdier
ikke lader Tid og Kraft tilbage til at leve i Tanken om deres
Bestaaen. Den etiske Religionsfilosofi danner en Parallel til den
erkendelsesteoretiske. I sine klassiske Tider bestemmer Religionen
direkte og indirekte Vurderingen af alle Handlinger og alle
Livsforhold. Det religiøse Problem rejser sig paa dette Punkt, naar
der har udviklet sig en selvstændig Etik, der ad sin egen Vej søger
Begrundelse og Motivering for Vurderingen af de menneskelige
Handlinger. Og paa dette Punkt stilles Problemet paa den skarpeste og
mest indtrængende Maade, ti den sluttelige Maalestok for Bedømmelsen
af Religionen maa høre hjemme paa det etiske Omraade, hvor det sidste
Opgør af det aandelige Regnskab finder Sted. ---

Indenfor den her angivne Ramme vilde vi kunne finde Plads for alle
væsentlige Sider ved det religiøse Problem. Jeg vil behandle dette
Problem, ikke blot med den intellektuelle Interesse, der vækkes ved
et stort og omfattende Emne, men ogsaa i den Stemning, der
fremkommer ved Bevidstheden om at staa overfor en aandelig Livsform,
i hvilken Menneskeslægten gennem Aartusinder har nedlagt nogle af
sine dybeste og inderligste Erfaringer.

\medskip

\noindent\textbf{3.} Inden vi gaa ind paa de enkelte Dele af
Religionsfilosofien, vil det være hensigtsmæssigt at give en noget
nærmere Karakteristik af det religiøse Problems Indhold.

I den protestantiske Verden er det temmeligt almindeligt anerkendt, at
der maa skelnes mellem Viden og Tro. Den katolske Kirke anerkender
ikke denne Distinktion i dens fulde Udstrækning og anerkender derfor
heller ikke noget religiøst Problem paa det intellektuelle Omraade.
Den hævder, at Guds Tilværelse kan bevises ad de Veje, Thomas
Aqvinas, støttet til Aristoteles, paaviste, og den har endog erklæret
det for en Trossag at antage, at disse Beviser slaa til. Traditionen
fra Middelalderens store Tid, da ``man ikke tvivlede og derfor var
nærmere ved Sandheden'' (for at bruge en Sætning af én af Thomas
Aqvinas' moderne Biografer), holdes stadigt i Hævd, og det
anerkendes ikke, at der virkeligt er kommen en Disharmoni ind i det
aandelige Liv. Distinktionen mellem Viden og Tro indrømmer, at en
saadan Disharmoni er til Stede, --- at Aandslivets Enhed er brudt. Er
nu denne Distinktion selv kun en Station paa Vejen, idet Udviklingen
vil føre til en Udskillelse af al Religion af Aandslivet, fordi kun
det kan bestaa, der kan træde i harmonisk Forhold til Livets andre
Elementer? --- Hvad enten man nu antager, at det kun er Religionens
klassiske Tid, der er forbi, eller at dens Tid overhovedet er forbi,
saa kan det Spørgsmaal ikke skydes til Side, om der gaar aandelig
Værdi tabt ved den Opløsningsproces, der er indtraadt, eller om der
kan paavises Erstatning for den aandelige Enhed og Sluttethed, som
tabtes. Allerede denne Betragtning viser, at Spørgsmaalet om
Værdiens Bestaaen er en væsentlig Side ved det religiøse Problem,
hvilken Løsning af dette man saa end helder til. Det er ikke mindst
fra denne Side, at Filosofien i det sidste halvandet Aarhundrede (fra
Rousseau's, Lessing's og Kant's Tid af) har behandlet det religiøse
Problem. Baade Romantikere og Positivister have været opmærksomme paa
denne Side ved det. --- Det maa jo vel erindres, at selvom man kunde
paavise, at der ingen aandelig Energi gik tabt ved Religionens
Opløsning, var det ikke sagt, at Energien under de nye Forhold fik en
ligesaa værdifuld Anvendelse som i Religionens klassiske Tider.
Ligesaa lidt paa det aandelige som paa det materielle Omraade er
Energiens Bestaaen uden videre det Samme som Værdiens Bestaaen.
Energien kan bestaa, uden at Værdien bestaar, medens det Omvendte
ikke vel kan tænkes, da Energi i og for sig vil være en Betingelse
for det Værdifuldes Bestaaen og Udvikling. Mange Fritænkere gaa uden
videre ud fra, at det menneskelige Liv vil blive kraftigere og
rigere, naar Religionen falder bort, --- en Antagelse, som ikke er
selvfølgelig, og som hviler paa Forudsætningen om, at der er psykiske
Ækvivalenter til Stede, --- Ækvivalenter baade hvad Værdi og hvad
Energi angaar. Det vil da blive Opgaven at paavise disse Ækvivalenter
og saaledes godtgøre Værdiens Bestaaen. Men det er et stort
Spørgsmaal og en væsentlig Side ved det religiøse Problem, om
saadanne Ækvivalenter kunne paavises. --- Iøvrigt vil Problemet om
Værdiens Bestaaen ogsaa opstaa indenfor Kirken, ti Religionen og det
religiøse Liv undergaar stadige Forandringer; og Spørgsmaalet bliver,
om Værdierne fra Religionens klassiske Tider holde sig gennem disse
Forandringer, f.~Ex., om Urkristendommen, baade hvad Værdi og hvad
Energi angaar, er bevaret i vore Dages Kristendom. I sin skarpeste
Form er dette Problem fremtraadt i den Paastand: ``det nye
Testamentes Kristendom er ikke til.'' ---

Men det religiøse Problem angaar ikke blot Værdiens Bestaaen indenfor
den menneskelige Verden. Det stiller ikke blot et psykologisk, men
ogsaa et kosmologisk Spørgsmaal, nemlig om Forholdet mellem det, der
for Mennesket staar som det højeste Værdifulde, og den hele
Tilværelse. Er der nogen Sammenhæng mellem Værdierne og Tilværelsens
Love og Kræfter? Ere maaske endog disse Love og Kræfter tilsidst
bestemte ved de højeste Værdier? Eller kunne vi maaske slet ikke
tillægge Værdibegrebet nogen Gyldighed ud over det menneskelige Livs
Omraade?

Den Hypotese, som senere (i den psykologiske Religionsfilosofi) skal
søges begrundet, gaar ud paa, at Sætningen om Værdiens Bestaaen er
det egentlige religiøse Axiom, der i de forskellige Religioner og paa
de forskellige religiøse Standpunkter udtrykkes paa forskellig Maade.
Det religiøse Problem angaar altsaa det Spørgsmaal, hvorvidt dette
Axiom virkeligt kan antages at gælde. Og tillige vil dette Axiom,
saafremt det udtrykker al Religions Grundtanke, kunne bruges som en
Maalestok for den enkelte Religions eller det enkelte religiøse
Standpunkts Konsekvens og Betydning. Endeligt vil --- som allerede
bemærket --- Forholdet mellem Etik og Religion ved Hjælp af dette
Axiom kunne udtrykkes paa en meget simpel Maade: hvorledes forholder
Troen paa Værdiens Bestaaen sig til Arbejdet paa at finde, frembringe
og bevare Værdier? ---

Dersom Sætningen om Værdiens Bestaaen skulde drøftes i hele dens
Rækkevidde, vilde Diskussionen blive uafsluttelig, og der vilde
desuden snart mangle klare Udgangspunkter for at føre den. Jeg
fremhæver derfor straks de Grænser, indenfor hvilke min Undersøgelse
vil bevæge sig, og som maa haves i Erindring i alt Følgende. --- Naar
jeg forsøger paa at vise, at den nævnte Sætning er Grundtanken i al
Religion, vil dette Forsøg ikke være mislykket, fordi det viser sig,
at ingen Religion udtrykker Sætningen med Klarhed og Konsekvens. Det
vil være tilstrækkeligt, om der kan paavises en udtrykkelig Trang
eller Tendens til at fastholde denne Grundtanke, saaledes at
Maalestokken ved Vurderingen af Religionerne i deres indbyrdes
Forhold bliver den Grad, i hvilken de formaa at udtale og gennemføre
Sætningen. --- Naar der tales om Værdiens Bestaaen, saa kunde dette
forstaas saaledes, at Værdien ikke forsvinder, men at der stadigt er
Værdi i Tilværelsen, hvad enten den aftager eller voxer eller skifter
mellem Aftagen og Tiltagen. Men jeg opfatter Udtrykket Værdiens
Bestaaen i Analogi med Udtrykket Energiens Bestaaen, saaledes at
Sætningen hævder en stadig Bestaaen af Værdien gennem alle
Formforandringer. Det vil maaske herved blive nødvendigt at skelne
mellem potentiel og aktuel Værdi, ligesom man skelner mellem potentiel
og aktuel Energi, og den ene Distinktion er ligesaa tydelig og
ligesaa berettiget som den anden. --- Men Værdiens Bestaaen gennem
alle Formforandringer og trods Forskellen mellem Mulighed og
Virkelighed er dog kun et Minimum. Der vil kunne paastaas, at den
blotte Bestaaen af Værdien ikke er tilfredsstillende og egentligt
indeholder en Modsigelse, da Værdien kun kan bestaa ved at vokse, da
blot Gentagelse sløver, medens paa den anden Side Forandring kun er
værdifuld, naar den bringer Forhøjelse. At en Forøgelse af Værdien
ikke strider mod Energiens Bestaaen, behøver ingen nærmere
Paavisning; Værdiforøgelsen forudsætter blot en ny og fuldkomnere
Anvendelse af Energien, ikke en Forøgelse af denne i det Hele. Det
er for Simpelhedens Skyld, at jeg begrænser Undersøgelsen til
Værdiens Bestaaen. Og det vil tillige vise sig, at den religiøse
Bevidsthed i Grunden ogsaa har begrænset sig hertil. Værdiens
Bestaaen vil stille os Vanskeligheder og Opgaver nok, selvom vi ikke
indlade os paa Alt, hvad der ligger i den uundgaaelige Konsekvens, at
Værdien kun kan bestaa, dersom den forøges. Dog maa vi ikke glemme,
at denne Konsekvens er til Stede. --- Endeligt kunde man, naar der er
Tale om Værdiens Bestaaen, spørge: hvilken Værdi og hvor megen Værdi?
Paa disse Spørgsmaal kommer jeg kun ind, forsaavidt Forskellen mellem
de forskellige Religioner og religiøse Standpunkter væsentligt beror
paa de forskellige Svar, de give paa dem. Det er klart, at Sætningen
om Værdiens Bestaaen kan antages, hvilke Svar man saa giver.
Hypotesen om, at Værdiens Bestaaen er den religiøse Grundsætning,
rokkes ikke ved, at de forskellige Religioner og religiøse
Standpunkter afvige i høj Grad fra hverandre i Henseende til
Beskaffenheden og Omfanget af det Værdifulde, paa hvis Bestaaen de
tro.

Værdi er den Egenskab ved Noget, at det medfører en virkelig
Tilfredsstillelse eller kan blive Midlet til en saadan. Værdien kan
altsaa være umiddelbar eller middelbar. Den umiddelbare Værdi søge vi
at fastholde, naar den er given, og at frembringe, naar den ikke er
given. Vi gøre den altsaa til Formaal for os. At vi gøre Noget til
Formaal, forudsætter omvendt, at vi have en Erfaring eller en
Formodning om dets umiddelbare Værdi. Middelbar Værdi har det, der
tjener til et Formaals --- altsaa en umiddelbar Værdis --- Opnaaelse.
Middelbar Værdi behøver ikke at være potentiel Værdi. Ti hvad der
tjener som Middel til at opnaa Noget, der har umiddelbar Værdi, faar
selv Værdi i vore Øjne, selvom dets Værdi staar som afledet. Der er
alle mulige Overgange mellem middelbar og umiddelbar Værdi, og ved
Motivforskydning gaar hin over til denne, saa at det, der først kun
havde Værdi som Middel, nu faar Værdi som Formaal. Den potentielle
Værdi derimod betyder ofte kun Muligheden af, at Noget faar Værdi som
Middel, --- altsaa Muligheden af middelbar Værdi. --- Hvilke de
Værdier ere, der erfares og derfor anerkendes paa de forskellige
Standpunkter, vil bero paa disse Standpunkters forskellige
Beskaffenhed. Det er et Væsens Natur, der bestemmer dets Trang, og
det er dets Trang, der bestemmer, hvad der faar Værdi for det. Det
religiøse Axiom viser altsaa Nødvendigheden af, at en Religions
Karakter bestemmes ved de Menneskers Natur og Trang, som hylde den.
Ti man kan ikke for Alvor tro paa Bestaaen af Værdier, som man ikke
tilnærmelsesvis kender af egen Erfaring. ---

Det er Kant's Filosofi, vi skylde Værdiproblemets Selvstændighed
overfor Erkendelsesproblemet. Han har lært os at skelne mellem
Vurdering og Forklaring. At denne Forskel dog ikke kan være en
fuldstændig Modsætning, ses af, at det maa bero paa Tilværelsens Love
og Kræfter, om Noget kan faa Værdi for os, og om vi kunne gøre det
Værdifulde til et Formaal, der kan virkeliggøres. Det religiøse
Problem kan derfor ikke skilles fra de andre filosofiske Problemer,
selvom det er heldigst at behandle det for sig selv. Det frembyder en
tydelig Analogi med de andre Problemer. Det filosofiske Grundproblem
angaar Kontinuiteten i Tilværelsen i dens Forhold til de specielle og
individuelle Tilværelsesformer. Og dersom den Hypotese, at Sætningen
om Værdiens Bestaaen er det religiøse Axiom, viser sig at være
begrundet, saa angaar det religiøse Problem jo Kontinuiteten i
Tilværelsen fra et særeget Synspunkt og gaar dermed ind under
Grundproblemet som en særlig Form for dette.


%% ============================================================
%%  II. ERKENDELSESTEORETISK RELIGIONSFILOSOFI
%% ============================================================
\chapter{Erkendelsesteoretisk Religionsfilosofi}
\label{ch:erkendelsesteori}

\section{Forstaaelse}
\label{sec:forstaaelse}

\begin{flushright}
\textit{Ihm ziemt's die Welt im Innern zu bewegen.}\\[2pt]
\textsc{Goethe}.
\end{flushright}

\bigskip

\noindent\textbf{4.} Mod den Formodning, som jeg har fremført, at
Religionens Kærne er Tro paa Værdiens Bestaaen, kunde det synes at
stride, at teoretiske Motiver jo dog have spillet, og endnu hos Mange
spille en betydelig Rolle ved religiøse Forestillingers Udvikling. Jeg
har desuden allerede bemærket, at da Religionen i sine klassiske Tider
er Alt for Mennesket, fyldestgør den i saadanne Tider ogsaa dets
Erkendelsestrang, det vil sige, giver det Midler til og Former for
Forstaaelse og Forklaring af Tilværelsen. Der existerer i Religionens
klassiske Tider ingen selvstændig Erkendelsestrang, saalidet som
særlige Midler og Former til at tilfredsstille en saadan Trang.
Sondringen mellem Viden og Tro, mellem Forstaaelse og Vurdering
kommer først frem i de kritiske Tider, det vil sige, i de Tider, da
Religionen er bleven et Problem. Denne Sondring er først begyndt at
trænge igennem efter adskillige voldsomme Sammenstød mellem Religion
og Videnskab. Kun modstræbende er man fra religiøs Side begyndt at
lære, at det ikke er Religionens Sag at give os en videnskabelig
Forstaaelse af Verden. Hvad nu alle Teologer gentage, at Bibelen ikke
skal lære os Naturvidenskab, vilde Ingen høre, da Bruno, Galilei og
Spinoza sagde det. Kætterforfølgelsen havde sit Forløb, og først
bagefter saa man Rigtigheden af, hvad Kætterne havde sagt. Paa andre
Omraader gentager den samme Strid sig; i vore Dage er det
Aandsvidenskaben, hvis Selvstændighed det gælder.

Enhver stor Religion optræder i Historien med en Verdensanskuelse,
idet den paa ejendommelig Maade sammenfatter i sig Tidens hele
Forestillingskreds. Kristendommen f.~Ex.\ stod i de første
Aarhundreder ogsaa ved sit Tankeindhold som tiltrækkende for Mange.
Tatianos gik saaledes over til Kristendommen, fordi denne
``Barbarernes Filosofi'', som han kaldte den, syntes ham at forklare,
hvad han før ikke havde forstaaet: ``Verdens Tilblivelse''. De gamle
Apologeter stode overhovedet paa det Standpunkt, at Kristendommen
indeholdt den bedste Løsning af Tankens Gaader. Naar Platon og andre
græske Filosofer syntes at antyde lignende Løsninger, saa var det
ifølge Justinus Martyr enten at forklare ved, at de vare blevne
paavirkede af de israelitiske Profeter, eller ogsaa ved, at ``Ordet''
(Logos), før det blev Kød, havde virket oplysende i Menneskenes Indre.
I Sammenligning med de tidligere Religioner havde Kristendommen ---
som det overhovedet er Tilfældet med de højere Religioner i Forhold
til de lavere --- en vis rationel Karakter og især en Simpelhed, en
enkel Storhed, der maatte virke tiltrækkende paa Mange, som kun
kendte de tidligere Religioner med deres brogede Mytologier.
Forsaavidt den Verdensanskuelse, der var given i de bibelske Skrifter,
ikke var udførlig nok, supplerede Kirkefædre og Skolastikere den ved
den græske Videnskabs Hjælp.

Ligesom Kristendommen for den vestlige Menneskehed staar Buddhismen
for den østlige Menneskehed ikke blot som en stor Religion, men som en
stor Verdensanskuelse, hvis Ideer have en vis rationel Karakter og
ikke kunne forstaas uden Kendskab til den forudgaaende indiske
Tænkning.

Naar der er Tale om Modsætningen mellem religiøs og videnskabelig
Forstaaelse af Tilværelsen, lægger man i Regelen Vægten paa, at
Videnskaben --- især paa Astronomiens, Geologiens og Biologiens
Omraade --- har ført til Resultater, som stride mod Religionens
overleverede Lærdomme. En stor Del af de moderne Apologeters Arbejde
gaar da ogsaa ud paa at vise Foreneligheden af de overleverede
religiøse Forestillinger med de videnskabelige Resultater, der
virkeligt staa fast. Man anvender overfor den moderne Naturvidenskab
med større eller mindre Held --- og med større eller mindre Smag ---
den samme Tillempelsesproces, som Oldtidens Apologeter anvendte
overfor den græske Videnskab. Men til Belysning af det religiøse
Problem i dets Helhed er hele denne Diskussion for og imod af ringe
Interesse, uden maaske forsaavidt den afgiver karakteristiske
Exempler paa den Tillempelse, det religiøse Forestillingsliv ---
saavel som alt Levende --- undergaar eller stræber at undergaa overfor
nye Livsforhold.

Hvad der derimod er Hovedsagen, er, at det er to helt forskellige
Arter eller Typer for ``Forstaaelse'', der her staa overfor hinanden.
Det er to Maader at tænke paa og at forklare Begivenheder paa, ---
to Maader, der ere saa forskellige, at naar man engang er begyndt med
den ene, vil man ikke kunne finde noget Punkt, hvor det er berettiget
og muligt at gaa over til den anden. Den hele aandelige Tilstand er
saa forskelligartet, at Ordet ``Forstaaelse'' i Virkeligheden her
bruges i to helt forskellige Betydninger. Lad mig tage et enkelt
Exempel. En pludselig Storm og Springflod, der volder mange Fiskeres
Undergang paa Havet, forklares naturvidenskabeligt ved Luftens og
Havets Tilstande i de foregaaende Dage. Men i Ligtalen over de
Forulykkede findes ``Forstaaelsen'' af hvad der er sket i, at Gud
vilde give de Efterlevende et Tegn, for at de skulde omvende sig. Paa
hvilket Punkt i Naturaarsagernes Række kan man nu tænke sig, at denne
guddommelige Indgriben har fundet Sted? Naturaarsagerne danne en
sammenhængende Række, i hvilken intet Led kan tages ud. Der maa
vælges mellem de to ``Forstaaelser''. Den ene Art Forstaaelse gaar ud
paa saa vidt muligt at danne sammenhængende Rækker af Led, der hvert
for sig kunne paavises i Erfaringen. Jo mere Luftens og Havets
Tilstand i Katastrofens Øjeblik er bragt i sammenhængende Rækkefølge
med deres Tilstande i de foregaaende Dage, des bedre siger man ---
efter den ene Type for Forstaaelse --- at man forstaar Katastrofen.
Efter den anden Type forstaas denne ved den Værdi, den faar
formedelst sin Indflydelse paa deres Sind, som lide under den, og
denne Værdi betragtes som tilsigtet, --- som Formaalet for en
guddommelig Indgriben.

Den videnskabelige Type for Forstaaelse har især udviklet sig paa
klar og frugtbar Maade i Løbet af de sidste tre Aarhundreder. Den
udelukker ikke en religiøs Vurdering, men den udelukker en
Forvexling af denne Vurdering med egentlig Forklaring. Det er den
filosofiske Erkendelsesteoris Opgave at undersøge de Forudsætninger
og de Tankeformer, som ligge til Grund for videnskabelig Erkendelse,
og det vil være af religionsfilosofisk Interesse at fremdrage disse
Forudsætninger og Tankeformer og sammenholde dem med Grundlaget for
den religiøse ``Forstaaelse''. Derigennem ville vi kunne faa
Oplysning om, hvorvidt Religionen fremdeles vil kunne have Betydning
for Tilfredsstillelsen af vor Erkendelsestrang.

\subsection{Aarsagsforklaring}
\label{ssec:aarsag}

\noindent\textbf{5.} At forstaa er at føre et hidtil Ubekendt
tilbage til noget Bekendt. At forstaa Noget er ikke altid det Samme
som at paavise dets Aarsag; man kan forstaa, hvad Noget er, uden at
forstaa, hvorfor det er det eller idethele er til. Og selve
Forstaaelsen af, hvorfor Noget er som det er, kan være af forskellig
Art.

Forstaaelse kan betyde Genkendelse (Perception). Jeg forstaar et
Sprog, naar jeg kender Lydene, saa at jeg kan forbinde den rette
Betydning med dem. Jeg forstaar, hvad det Hvide mellem Buskene er,
naar jeg bliver klar over, at det er en Taagestribe, som Maanen
skinner paa. Jeg forstaar en Sætning, naar den Forbindelse mellem
Subjekt og Prædikat, den Talende eller Skrivende har villet hævde,
genkendes af mig. --- I disse Exempler beror Forstaaelsen paa
Opdagelse af Identitet.

Til Forstaaelse af en Sætning kan der dog fordres noget Mere. Fordi
jeg forstaar, hvad Meningen er med den enkelte Sætning, forstaar jeg
ikke, hvorfor Sætningen er bleven fremsat. Denne sidste Forstaaelse
kommer først, naar jeg kan udlede Sætningen af andre Sætninger,
indtil jeg kommer til Sætninger, der forud staa fast for mig.
Forstaaelsen betyder her Begrundelse. Sætningen ``A er C'' forstaar
jeg ikke (i denne Betydning af Ordet), før jeg ser, at den følger
af, at ``A er B'' og ``B er C''. --- Her beror Forstaaelsen paa
Rationalitet.

En Begivenhed i den sjælelige eller i den materielle Verden forstaar
jeg (i første Betydning), naar jeg véd, hvad der er sket. Og jeg
forstaar den (i anden Betydning), naar jeg véd, hvorfor jeg antager,
at det er sket. Men der er endnu en tredie Art af Forstaaelse mulig.
At forstaa en Begivenhed kan betyde at udlede den af tidligere
Begivenheder, saa at der af disses Indtræden kan sluttes til dens
Indtræden. Her betyder Forstaaelsen Aarsagsforklaring og bestaar i
Opdagelsen af Kausalitet.

Identitet, Rationalitet og Kausalitet ere beslægtede. --- Rationalitet,
Forholdet mellem Grund (Præmisser) og Følge (Konklusion),
forudsætter Identitet. Sætningen ``A er C'' er paa en Maade allerede
indeholdt i de to Sætninger ``A er B'' og ``B er C'', og er
egentligt kun et andet Udtryk for disse to Sætningers Indhold, et
Udtryk, som fremkommer ved, at jeg i Stedet for B i den første
Sætning sætter B's Prædikat (C) i den anden Sætning, hvilket jeg maa
gøre, naar de to B'er virkeligt ere identiske. Konklusionen
fremkommer ved en Kombination af Præmisserne. --- Kausaliteten
forudsætter igen Rationaliteten. Forholdet mellem Aarsag og Virkning
er analogt det mellem Grund og Følge. To Begivenheder staa i
Aarsagsforhold til hinanden, naar de ikke blot følge efter hinanden i
Tiden, men naar den, der indtræder først, tillige forholder sig til
den anden som Grund til Følge. Kausaliteten indeholder Rationalitet,
men forudsætter desuden, at der er sket en Forandring, en Overgang i
Tiden fra en Begivenhed til en anden. Virkningen ligger i en vis
Forstand i Aarsagen, men er dog Mere end en rent logisk Omsætning af
denne i en anden Form. Det er ved Overgangen fra Aarsag til Virkning
ikke blot vor Tanke, der gaar over fra en Form til en anden, men der
foregaar en Overgangsproces i selve Virkeligheden. Det store
Spørgsmaal er, om alle Forhold mellem virkelige Begivenheder ere
Aarsagsforhold.

Naar man opstiller den Sætning, at enhver Begivenhed har sin Aarsag,
--- den saakaldte Aarsagssætning, --- saa slaar man derved det
Princip fast, at Tilværelsen er forstaaelig i Ordets tredie
Betydning. Denne Sætning kan aldrig blive Mere end en Hypotese, da
vi aldrig kunne forstaa Alt, paavise Aarsager til alle Begivenheder.
Men den indeholder den Forudsætning, med hvilken vi uvilkaarligt møde
overfor Begivenheder, der vække vor Opmærksomhed; i hvert Tilfælde
er det den Forudsætning, som bringer os til at se os om efter
tidligere Begivenheder, naar en Begivenhed er indtraadt.

Vi kunne her ganske se bort fra Spørgsmaalet om Aarsagssætningens
Begrundelse og Oprindelse, da de indbyrdes afvigende Opfattelser,
der her kunne gives, ikke behøve at faa religionsfilosofisk
Betydning. Hvad enten Anerkendelsen af denne Sætning skyldes Vane
eller en indre Nødvendighed i vor Natur, saa rører Trangen til at
finde Aarsager sig mere eller mindre hos Alle, og den religiøse
Opfattelse bærer saavel som den videnskabelige Vidnesbyrd derom. Ja,
de religiøse Forestillinger have netop en særegen erkendelsesteoretisk
Interesse derved, at de vidne om, hvor dybt Trangen til at finde
Aarsager ligger i den menneskelige Natur, idet den ytrer sig selv
dér, hvor endnu ingen selvstændig videnskabelig Interesse er
opstaaet. Modsætningen mellem Religion og Videnskab kommer først
frem, naar der ses hen til Beskaffenheden af de Aarsager, man
anerkender som rette, virkelige Aarsager (\textit{veræ causæ}). Baade
religiøs og videnskabelig Forstaaelse er en Tilbageføren af det
Ubekendte til det Bekendte; Forskellen mellem dem beror paa, hvad
dette Bekendte er.

Den videnskabelige Forstaaelse kræver, at Aarsagen til en Begivenhed
skal paavises i Begivenheder, der selv ligesaa vel godtgøre sig i
Erfaringen, som den Begivenhed, der skal forklares. Opgaven bliver
her at fortolke Begivenhedernes Række ved at finde saa mange enkelte
Led i denne Række og saa nøje Sammenhæng mellem dem som muligt. Det
gælder om at fortolke Naturen ved Naturen selv, ligesom man fortolker
et Sted i en Bog ved at sammenholde det med andre Steder i samme Bog.
Den religiøse Forstaaelse stiller ikke det Krav, at Aarsagen skal
staa i kontinuerlig Sammenhæng med Virkningen og selv paavises som en
Erfaring ligesaa vel som den Begivenhed, der skal forklares. Den
undersøger ikke sine Aarsager paa den kritiske Maade, som
Videnskaben gør det. Den forklarer ikke Naturen ved Naturen selv, men
ved Noget, der er forskelligt fra eller staar udenfor Naturen,
``udefra støder'' til denne. Saaledes optræder i hvert Tilfælde den
religiøse Forstaaelse hyppigst i Historien; om det er nødvendigt, at
den optræder i denne Form og derfor i udpræget Modsætning til den
videnskabelige Forstaaelse, vil senere vise sig. I Modsætning til
denne religiøse Forstaaelse kan den specielle Form, Erfaringsvidenskaben
--- baade paa det sjælelige og det materielle Omraade --- giver
Aarsagssætningen, betegnes som de naturlige Aarsagers Princip. Hvis
det ikke kan anvendes, forstaa vi Intet, naar det er den
videnskabelige Type for Forstaaelsen og den tilsvarende intellektuelle
Habitus, der raader hos os.

\medskip

\noindent\textbf{6.} Men hvorpaa grunder sig da de naturlige Aarsagers
Princip? --- Der kan paavises en hel Række Grunde til, at dette
Princip har arbejdet sig frem og mere og mere bestemmer den Maade,
paa hvilken den menneskelige Bevidsthed stiller sig overfor
Begivenhederne, naar det gælder om at forstaa dem.

Det er ikke nok at udtænke en Aarsagsforklaring. Det maa ogsaa
paavises, at denne Forklaring bekræftes ved Erfaring. En saadan
Bekræftelse (Verifikation) opnaas, naar vi se, at ved samme Forholds
Genindtræden indtræder samme Begivenhed som forrige Gang. Men det er
klart, at en saadan Bekræftelse kun er mulig, naar Aarsagen selv er
en Begivenhed, der kan fremtræde i utvivlsom Erfaring. Hvis f.~Ex.\
den ovenfor (4) omtalte Forklaring af Stormen og Springfloden er
rigtig, maa disse Begivenheder indtræde igen, naar de samme Tilstande
i Luft og Hav ere til Stede, og dette maa kunne godtgøres. At udlede
Begivenhederne af guddommelige Hensigter er derimod ingen
videnskabelig Forklaring; ti om saadanne Hensigter i det enkelte
Tilfælde ere til Stede, kan ikke godtgøres, og Forklaringen kan
derfor ikke prøves. Det var netop af denne Grund, ikke af Uvillie
mod Religionen, at Galilei, Bacon og Descartes erklærede de saakaldte
Hensigtsaarsager for uvidenskabelige.

Men dertil kommer, at kun de naturlige Aarsagers Princip
tilfredsstiller Trangen til speciel Forklaring af den enkelte,
bestemte Begivenhed. Galilei viste saaledes, at det i Debatten mellem
de forskellige stridende astronomiske Hypotheser ikke kunde nytte at
beraabe sig paa Guds Indgriben, da det for Guds Almagt er ligesaa
let at lade Solen bevæge sig om Jorden som at lade Jorden bevæge sig
om Solen. Denne Aarsag kan altsaa ikke hjælpe os til et Valg mellem
Hypotheserne. Den er en Nøgle, der lukker op overalt; men det er
ikke en saadan Hovednøgle, Videnskaben søger; Videnskabens Interesse
er det at undersøge de forskellige Laases særegne Beskaffenhed, og
den fordrer derfor en egen Nøgle til enhver Laas. Dersom en
Hovednøgle var tilstrækkelig, vilde Laasenes Forskellighed blive
betydningsløs --- og indeholde et uløst Problem! --- At beraabe sig
paa Guds Villie giver ingen videnskabelig Forstaaelse, men er ---
videnskabeligt set --- en Indrømmelse af, at man ikke forstaar
Begivenheden.

Fremdeles. I Tvivlstilfælde have vi kun ét Kriterium, ved hvilket vi
kunne skelne mellem Indbildning og Virkelighed, nemlig den faste,
urokkelige Sammenhæng mellem det enkelte Fænomen og vor hele øvrige
Erfaring. Jo mere løsrevet det enkelte Fænomen staar i Forhold til
vor øvrige Erfaring, --- jo større Spring der er mellem Leddene i
Begivenhedernes Række, --- jo mere uensartede Leddene ere, --- des
mindre Sikkerhed have vi for at staa overfor en Virkelighed. I
Tvivlstilfælde vedblive vi derfor at søge efter Mellemled i
Erfaringen, indtil vi have fundet en kontinuerlig og urokkelig
Sammenhæng. Det er klart, at Anvendelsen af dette
Virkelighedskriterium er Et med Anvendelsen af de naturlige Aarsagers
Princip.

Den dybeste Begrundelse for de naturlige Aarsagers Princip finder jeg
i en Trang til Kontinuitet, som ligger i vor Bevidstheds Natur, og
som allerede lægger sig for Dagen i den almindelige Trang til at
finde Aarsager. Vor Bevidsthed stræber uvilkaarligt at sammenfatte og
sammenholde sine Elementer og at bevare Enheden med sig selv, ogsaa
naar den staar overfor et nyt Indhold. Det nye Indhold stræber den da
at forme, ændre og lægge til Rette, saaledes at det Nye viser sig at
være en ny Form for et Gammelt. Selvgenkendelse bliver da mulig trods
Indholdets Skiften, og det personlige Livs Enhed afspejler sig i den
Kontinuitet, som derved tilvejebringes. Erkendelsesarbejdet hænger
her nøje sammen med Arbejdet paa Karakterens Udvikling. Om en
Karakter tale vi --- i udpræget Forstand --- kun, hvor der hersker
Enhed og Kontinuitet i den hele Stræben og ikke atter og atter
springes over til noget Nyt. I Analogi hermed søger Erkendelsen at
finde saa megen Enhed og Kontinuitet som muligt, at paavise det Nye
som Omsætning eller Fortsættelse af det Gamle, og kun hvor dette
lykkes, er en fuld Tilegnelse og Forstaaelse vunden. De naturlige
Aarsagers Princip, der fører til at bygge saa lange og saa
kontinuerlige Rækker som muligt, er derfor ingen Tilfældighed, men
hænger nøje sammen med Personlighedens eget Væsen. Det er derfor
urigtigt, naar Personlighed og Videnskab sættes i skarp Modsætning
til hinanden. Videnskabeligt Arbejde er et Personlighedsarbejde. Og
jo mere dette Arbejde lykkes, des mere nærmer Aarsagsforklaringen sig
til Begrundelsen og Genkendelsen. Forskellighederne i Tilværelsen ere
reducerede til det mindst Mulige, naar vi med ét Blik, ``under
Evighedens Synspunkt'' overskue store Egne af Tilværelsen.

\subsection{Rummets Verden}
\label{ssec:rummet}

% [text to be added]

\subsection{Tidsløbet}
\label{ssec:tidsloeb}

% [text to be added]

\section{Afsluttende Tanker}
\label{sec:afsluttende}

% [text to be added]

\section{Tanke og Billede}
\label{sec:tanke-billede}

% [text to be added]

%% ============================================================
%%  III. PSYKOLOGISK RELIGIONSFILOSOFI
%% ============================================================
\chapter{Psykologisk Religionsfilosofi}
\label{ch:psykologi}

\section{Religiøs Erfaring og religiøs Tro}
\label{sec:erfaring-tro}

\subsection{Religiøs Erfaring}
\label{ssec:erfaring}

% [text to be added]

\subsection{Religiøs Tro}
\label{ssec:tro}

% [text to be added]

\section{Religiøse Forestillingers Udvikling}
\label{sec:forestillinger}

\subsection{Religion som Drift}
\label{ssec:drift}

% [text to be added]

\subsection{Polyteisme og Monoteisme}
\label{ssec:polyteisme}

% [text to be added]

\subsection{Den religiøse Erfaring og Traditionen}
\label{ssec:tradition}

% [text to be added]

\subsection{Religionspsykologiens videnskabelige Afslutning}
\label{ssec:afslutning-psyk}

% [text to be added]

\section{Dogmer og Symboler}
\label{sec:dogmer}

% [text to be added]

\section{Sætningen om Værdiens Bestaaen}
\label{sec:vaerdien}

\subsection{Nærmere Bestemmelse af Sætningen om Værdiens Bestaaen
  og af dens Forhold til Erfaringen}
\label{ssec:vaerdi-bestemmelse}

% [text to be added]

\subsection{Psykologisk-historisk Drøftelse af Sætningen om
  Værdiens Bestaaen}
\label{ssec:vaerdi-psykologisk}

% [text to be added]

\subsection{Almindelig filosofisk Drøftelse af Sætningen om
  Værdiens Bestaaen}
\label{ssec:vaerdi-filosofisk}

% [text to be added]

\section{Personlighedsprincipet}
\label{sec:personlighed}

\subsection{Personlighedsprincipets Betydning og Berettigelse}
\label{ssec:pers-betydning}

% [text to be added]

\subsection{Hovedgrupper af personlige Forskelligheder}
\label{ssec:pers-grupper}

% [text to be added]

\subsection{Buddha og Jesus}
\label{ssec:buddha-jesus}

% [text to be added]

\subsection{Er Personlighedsprincipet et Princip for Væxt
  eller for Opløsning?}
\label{ssec:pers-vaext}

% [text to be added]

\subsection{Læg og Lærd}
\label{ssec:laeg-laerd}

% [text to be added]

%% ============================================================
%%  IV. ETISK RELIGIONSFILOSOFI
%% ============================================================
\chapter{Etisk Religionsfilosofi}
\label{ch:etik}

\section{Religion som Grundlag for Etik}
\label{sec:grundlag-etik}

% [text to be added]

\section{Religion som Form for aandelig Kultur}
\label{sec:aandelig-kultur}

\subsection{Psykologisk Betragtning}
\label{ssec:psyk-betragtning}

% [text to be added]

\subsection{Sociologisk Betragtning}
\label{ssec:sociol-betragtning}

% [text to be added]

\section{Urkristendom og moderne Kristendom}
\label{sec:urkristendom}

% [text to be added]

\section{Vi leve paa Virkeligheder}
\label{sec:virkeligheder}

% [text to be added]

%% ============================================================
%%  NOTER
%% ============================================================
\chapter*{Noter}
\addcontentsline{toc}{chapter}{Noter}
\label{ch:noter}

% [notes to be added]

\end{document}
